\uuid{DV6i}
\titre{Pgcd}
\auteur{bodin ; croizat ; sacre}
\organisation{exo7}
\licence{CC BY-NC-SA 4.0}
\datecreate{2022-03-08}
\niveau{L1}
\module{Algèbre}
\chapitre{Arithmétique dans Z}
\sousChapitre{Pgcd, ppcm, algorithme d'Euclide}
\difficulte{1}
\theme{pgcd}
\exo7qcmid{hilisit.arithmetique.004}

\contenu{
\begin{question}
\qclassification{103.03}
    Quel est le reste $r$ dans la division euclidienne de $145$ par $13$ ?
    \begin{answers}
        \bad{$r=0$}
        \good{$r=2$}
        \bad{$r=7$}
        \bad{$r=-11$}
    \end{answers}
    \begin{explanations}
     La division euclidienne de $145$ par $13$ nous donne l'écriture : $145 = 13 \times 11 + 2$. Le reste $r$ est donc $2$.\\
     Il est vrai que $145 = 13 \times 12 - 11$, mais cela ne correspond pas à une division euclidienne puisque le reste $r$ n'est pas compris entre $0$ et $13-1=12$.
    \end{explanations}
\end{question}
}
